\scanpage{53}
$S^{(\infty)}\longrightarrow G=\widehat S$, les hom
\inserted{\unsure{nuls} sur $S^0$} se recollent et définissent
un hom de \unsure{préschémas} formels
\[
S^{(\infty)}\longrightarrow G
\]
dont la connaissance \deleted{implique} \inserted{équivaut à}
la connaissance de l'ext. donnée $E$ de $M$ par $N$, non pas
sur $S$ lui-même, mais dans le voisinage inf. de
\unsure{la valeur} de la section de $S$ sur $R$, et
\unsure{partt} après changement de base.

\figure{\begin{tikzpicture}
\draw (0,0)--(.35,.45)--(.52,.4)--(2.7,.5)--(5.7,.42)--(9.7,.8);
\end{tikzpicture}}

La formation de $S^{(\infty)}\unsure{-|\to}\unsure{G|C}$
est compatible aux changements de base \deleted{\unsure{sur}}
$S'\longrightarrow S$. Il y a \unsure{vraisemblablement}
donc de la considérer, dans la mesure du possible, comme déduite
d'une situation ``\uuline{modulaire}''.

Supposons \deleted{\unsure{donc}} \deleted{\unsure{alors}} que l'on
ait \unsure{pris} $R=\operatorname{Spec}\Lambda$,
$\Lambda$ anneau local \inserted{noeth.} complet de car.
résiduelle $p$ (p.ex. un pro-anneau de Cohen).
de corps résiduel $k$, et qu'on se donne $M_0$ et $N_0$
sur $k$; On veut la variété modulaire
\inserted{formelle \unsure{(sur $R$ considéré comme schéma formel)}}
pour une extension fixée $\unsure{\xi_0}$ de $M_0$ par $N_0$.
Supposons $N_0$ de t.m.
Considérons aussi la variété modulaire formelle pour $\xi_0$,
l'extension triviale. \deleted{\unsure{On}} Alors:
\margin{63}
